\chapter{1977年黑龙江卷}

\section{解答题}

\begin{problem}
解答下列各题.

\begin{enumerate}
    \item 解方程：\(\sqrt{3x+4}=4\).

    \item 解不等式：\(\left|x\right|<5\).

    \item 已知正三角形的外接圆半径为 \(6\sqrt{3} \mathrm{~cm}\)，求它的边长.
\end{enumerate}
\end{problem}

\begin{solution}
\begin{enumerate}
    \item 根式有意义要求 \(3x+4\geqslant0\). 在此条件下，方程两边都是非负数，所以平方不改变方程的等价性：
    \[
    3x+4=4^2=16
    \]
    解得 \(x=4\). 代回原方程，左边为 \(\sqrt{3\times4+4}=\sqrt{16}=4\)，满足原方程，故解为 \(x=4\).
    \item 因为 \(5>0\)，绝对值不等式的定义给出
    \[
    \left|x\right|<5\Longleftrightarrow -5<x<5
    \]
    所以不等式的解集为 \(\left(-5,5\right)\).
    \item 设正三角形的边长为 \(a\)，高为 \(h\). 高平分底边，在所得直角三角形中由勾股定理得
\(h^2=a^2-\left(\frac a2\right)^2=\frac{3a^2}{4}\)，\(h=\frac{\sqrt3}{2}a\).
    正三角形的外心、重心和内心重合，且外心把中线按顶点到外心为中线长的 \(\frac23\) 分割，因此外接圆半径
    \[
    R=\frac23h=\frac a{\sqrt3}
    \]
    代入 \(R=6\sqrt3\mathrm{~cm}\)，得
    \[
    a=R\sqrt3=6\sqrt3\times\sqrt3=18\mathrm{~cm}
    \]
    所以正三角形的边长为 \(18\mathrm{~cm}\).
\end{enumerate}
\end{solution}

\begin{problem}
计算下列各题.

\begin{enumerate}
    \item \(\sqrt{m^2-2ma+a^2}\).

    \item \(\cos 78^\circ\cdot\cos 3^\circ+\cos 12^\circ\cdot\sin 3^\circ\).

    \item \(\arcsin\left(\cos\frac{\pi}{6}\right)\).
\end{enumerate}
\end{problem}

\begin{solution}
\begin{enumerate}
    \item 先把被开方数配方：
    \[
    m^2-2ma+a^2=(m-a)^2
    \]
    算术平方根取非负值，而 \(\sqrt{u^2}=\left|u\right|\)，所以
    \[
    \sqrt{m^2-2ma+a^2}=\sqrt{(m-a)^2}=\left|m-a\right|
    \]
    \item 利用余函数关系 \(\cos78^\circ=\sin12^\circ\)，原式化为
    \[
    \sin12^\circ\cos3^\circ+\cos12^\circ\sin3^\circ
    =\sin15^\circ
    \]
    再用 \(15^\circ=45^\circ-30^\circ\) 及正弦差角公式，得
    \[
    \sin15^\circ=\sin45^\circ\cos30^\circ-\cos45^\circ\sin30^\circ
    =\frac{\sqrt2}{2}\cdot\frac{\sqrt3}{2}-\frac{\sqrt2}{2}\cdot\frac12
    =\frac{\sqrt6-\sqrt2}{4}
    \]
    \item 先求括号内的值：
    \[
    \cos\frac\pi6=\frac{\sqrt3}{2}=\sin\frac\pi3
    \]
    反正弦函数的值域规定为 \(\left[-\frac\pi2,\frac\pi2\right]\)，而 \(\frac\pi3\) 在这个区间内，因此主值确定为
    \[
    \arcsin\left(\cos\frac\pi6\right)=\arcsin\left(\sin\frac\pi3\right)=\frac\pi3
    \]
\end{enumerate}
\end{solution}

\begin{problem}
解下列各题.

\begin{enumerate}
    \item 解方程：\(3^{x+1}-9^{\frac{x}{2}}=18\).

    \item 求数列 \(2\)，\(4\)，\(8\)，\(16\)，\(\cdots\) 前十项的和.
\end{enumerate}
\end{problem}

\begin{solution}
\begin{enumerate}
    \item 对任意实数 \(x\)，由于底数为正，指数运算满足
    \[
    9^{\frac x2}=\left(3^2\right)^{\frac x2}=3^x
    \]
    原方程化为
    \[
    3\cdot3^x-3^x=18\Longleftrightarrow 2\cdot3^x=18
    \Longleftrightarrow 3^x=9=3^2
    \]
    因为底数 \(3>0\) 且 \(3\ne1\)，指数函数 \(3^x\) 是一一对应的，所以 \(x=2\). 代回原方程可得 \(27-9=18\)，故解为 \(x=2\).
    \item 数列的相邻两项之比都是 \(2\)，所以它是首项 \(a_1=2\)、公比 \(q=2\) 的等比数列. 因为 \(q\ne1\)，前十项和为
    \[
    S_{10}=a_1\frac{q^{10}-1}{q-1}=2\frac{2^{10}-1}{2-1}=2(1024-1)=2046
    \]
    所以前十项的和为 \(2046\).
\end{enumerate}
\end{solution}

\begin{problem}
解下列各题.

\begin{enumerate}
    \item 圆锥的高为 \(6 \mathrm{~cm}\)，母线和底面半径成 \(30^\circ\) 角，求它的侧面积.

    \item 求过点 \(\left(1,4\right)\) 且与直线 \(2x-5y+3=0\) 垂直的直线方程.
\end{enumerate}
\end{problem}

\begin{solution}
\begin{enumerate}
    \item 设圆锥的母线长为 \(l\)，底面半径为 \(r\). 作圆锥的轴截面，得到一个直角三角形，其中高为 \(6\mathrm{~cm}\)，斜边为母线 \(l\)，且母线与底面半径的夹角为 \(30^\circ\). 因而
\(\sin30^\circ=\frac{6}{l}\)，\(l=\frac6{1/2}=12\mathrm{~cm}\).
    同一个直角三角形中，底面半径是夹角 \(30^\circ\) 的邻边，所以
    \[
    r=l\cos30^\circ=12\times\frac{\sqrt3}{2}=6\sqrt3\mathrm{~cm}
    \]
    圆锥侧面积等于底面圆周长乘母线的一半，即
    \[
    S_{\text{侧}}=\frac12(2\pi r)l=\pi rl
    =\pi\times6\sqrt3\times12=72\sqrt3\pi\mathrm{~cm}^2
    \]
    \item 将已知直线化为斜截式：
    \[
    2x-5y+3=0\Longleftrightarrow y=\frac25x+\frac35
    \]
    所以其斜率为 \(\frac25\). 垂直直线的斜率为负倒数 \(-\frac52\). 过点 \(\left(1,4\right)\) 的所求直线由点斜式给出
    \[
    y-4=-\frac52(x-1)
    \]
    两边乘以 \(2\) 并整理，得到
    \[
    2y-8=-5x+5\Longleftrightarrow 5x+2y-13=0
    \]
    因此所求直线方程为 \(5x+2y-13=0\).
\end{enumerate}
\end{solution}

\begin{problem}
如果 \( \triangle ABC \) 的 \( \angle A \) 的平分线交 \(BC\) 于 \(D\)，交它的外接圆于 \(E\)，那么 \( AB\cdot AC=AD\cdot AE \).

\bitmapfigure[width=4.8cm]{img/1977/heilongjiang/q5_fig1.png}
\end{problem}

\begin{solution}
因为角平分线交边 \(BC\) 于 \(D\)，又交外接圆于 \(E\)，所以 \(A\)，\(D\)，\(E\) 三点共线，且 \(D\) 在 \(BC\) 上. 由 \(AD\) 平分 \(\angle BAC\)，并注意到射线 \(AE\) 与射线 \(AD\) 重合，有
\[
\angle BAD=\angle DAC=\angle EAC
\]
另一方面，\(D\) 在 \(BC\) 上，所以 \(\angle ABD=\angle ABC\). 点 \(A\)，\(B\)，\(C\)，\(E\) 共圆，\(\angle ABC\) 与 \(\angle AEC\) 都是弦 \(AC\) 所对的圆周角，故
\[
\angle ABD=\angle AEC
\]
于是由两个角分别相等，得到
\[
\triangle ABD\sim\triangle AEC
\]
其中对应关系为 \(A\leftrightarrow A\)、\(B\leftrightarrow E\)、\(D\leftrightarrow C\). 因此对应边成比例：
\[
\frac{AD}{AC}=\frac{AB}{AE}
\]
交叉相乘即得
\[
AD\cdot AE=AB\cdot AC
\]
即 \(AB\cdot AC=AD\cdot AE\)，结论成立.
\end{solution}

\begin{problem}
前进大队响应毛主席关于“绿化祖国”的伟大号召，\(1975\) 年造林 \(200\text{ 亩}\)，又知 \(1975\) 年至 \(1977\) 年这三年内共造林 \(728\text{ 亩}\)，求后两年造林面积的年平均增长率是多少？
\end{problem}

\begin{solution}
设后两年造林面积的年平均增长率为 \(r\)，则每年相对于上一年的增长倍数为 \(q=1+r\). 由于造林面积为正，必须有 \(q>0\). 于是
\(S_{1975}=200\)，\(S_{1976}=200q\)，\(S_{1977}=200q^2\).
由三年总面积得

\[
200+200q+200q^2=728
\Longleftrightarrow q^2+q+1=\frac{728}{200}=\frac{91}{25}
\]

移项得 \(q^2+q-\frac{66}{25}=0\). 用求根公式解得


\[
q=\frac{-1\pm\sqrt{1+\frac{264}{25}}}{2}
=\frac{-1\pm\frac{17}{5}}{2}
\]

所以 \(q=\frac65\) 或 \(q=-\frac{11}{5}\). 由 \(q>0\) 舍去负值，故 \(q=\frac65\)，从而

\[
r=q-1=\frac65-1=\frac15=20\%
\]

所以后两年造林面积的年平均增长率为 \(20\%\).
\end{solution}

\begin{problem}
解方程：\(\lg\left(2^x+2x-16\right)=x\left(1-\lg 5\right)\).
\end{problem}

\begin{solution}
常用对数的真数必须为正，因此原方程的定义域先要求 \(2^x+2x-16>0\). 又因为 \(1=\lg10\)，所以右端可以化为

\[
x(1-\lg5)=x(\lg10-\lg5)=x\lg2=\lg(2^x)
\]

注意到 \(2^x>0\)，在原方程的定义域内两边都是有意义的常用对数，故原方程等价于
\[
\lg\left(2^x+2x-16\right)=\lg(2^x)
\]
常用对数函数在正数范围内是单调函数，从而真数相等：

\[
2^x+2x-16=2^x
\]

于是 \(2x-16=0\)，解得 \(x=8\). 检验定义域：
\[
2^8+2\times8-16=256>0
\]
所以 \(x=8\) 确实是原方程的解，且没有其他解.
\end{solution}

\begin{problem}
已知三角形的三边成等差数列，周长为 \(36 \mathrm{~cm}\)，面积为 \(54 \mathrm{~cm}^2\)，求三边的长.
\end{problem}

\begin{solution}
设三边按从小到大的顺序为 \(12-d\)、\(12\)、\(12+d\)，其中 \(d\geqslant0\). 这是因为三边和为 \(36\)，等差数列的中项等于平均数 \(36/3=12\). 三角形的最大边小于另外两边之和，故
\[
12+d<(12-d)+12\Longleftrightarrow d<6
\]
半周长为 \(s=18\). 根据海伦公式，

\[
54^2=18\left[18-(12-d)\right](18-12)\left[18-(12+d)\right]
 =18(6+d)\cdot6\cdot(6-d)=108(36-d^2)
\]

因此
\[
2916=108(36-d^2)\Longleftrightarrow d^2=9
\]
结合 \(d\geqslant0\)，得 \(d=3\). 所以三边分别为
\(12-3=9\mathrm{~cm}\)，\(12\mathrm{~cm}\)，\(12+3=15\mathrm{~cm}\).
最后，\(9+12>15\)，且代入海伦公式得到面积 \(54\mathrm{~cm}^2\)，满足题设，故答案成立.
\end{solution}

\section{附加题}

\begin{problem}
如图，\(AP\) 表示发动机的连杆，\(OA\) 表示它的曲柄. 当 \(A\) 在圆上作圆周运动时，\(P\) 在 \(x\) 轴上作直线运动，求 \(P\) 点的横坐标. 为什么当 \(\alpha\) 是直角时，\(\angle P\) 是最大？

\bitmapfigure[width=6.4cm]{img/1977/heilongjiang/q9_fig1.png}
\end{problem}

\begin{solution}
取 \(O\) 为原点，题图中的 \(x\) 轴为横轴. 设连杆长 \(AP=l\)，曲柄长 \(OA=R\)，并设 \(P\) 点横坐标为 \(p\)，于是 \(P=\left(p,0\right)\). 由 \(\angle AOP=\alpha\) 及 \(OA=R\)，点 \(A\) 的坐标为
\[
A=\left(R\cos\alpha,R\sin\alpha\right)
\]
由两点间距离公式，
\[
l^2=AP^2=\left(p-R\cos\alpha\right)^2+R^2\sin^2\alpha
\]
移项得
\[
\left(p-R\cos\alpha\right)^2=l^2-R^2\sin^2\alpha
\]
题图取的是连杆在曲柄右侧的分支，即 \(p-R\cos\alpha\geqslant0\)，所以
\[
p=R\cos\alpha+\sqrt{l^2-R^2\sin^2\alpha}
\]
这个位置存在的必要条件是 \(l^2-R^2\sin^2\alpha\geqslant0\)，因此上式就是 \(P\) 点的横坐标.

再说明角 \(P\) 的最大性. 设 \(\theta=\angle APO\)，这正是题目所记的 \(\angle P\). 在三角形 \(AOP\) 中，边 \(OA=R\) 对应角 \(\theta\)，边 \(AP=l\) 对应角 \(\alpha\)，由正弦定理得
\(\frac{\sin\theta}{R}=\frac{\sin\alpha}{l}\)，\(\sin\theta=\frac{R}{l}\sin\alpha\).
由所取分支 \(p-R\cos\alpha\geqslant0\). 在点 \(P\) 处用向量内积或余弦定理可得
\[
\cos\theta=\frac{p-R\cos\alpha}{l}\geqslant0
\]
所以 \(0\leqslant\theta\leqslant\frac\pi2\)，在此范围内 \(\theta\) 随 \(\sin\theta\) 单调增加. 当题设机构可以取到 \(\alpha=\frac\pi2\) 时，\(0\leqslant\alpha\leqslant\pi\) 范围内有 \(\sin\alpha\leqslant1\)，且等号仅在 \(\alpha=\frac\pi2\) 时成立. 因而 \(\sin\theta\)，以及 \(\theta=\angle P\)，在 \(\alpha\) 为直角时取得最大值.
\end{solution}

\section{参考题}

\begin{problem}
求曲线 \(y=\sin x\) 在 \(\left[0,\pi\right]\) 上的曲边梯形绕 \(x\) 轴旋转一周所形成的旋转体的体积.
\end{problem}

\begin{solution}
在 \(\left[0,\pi\right]\) 上有 \(\sin x\geqslant0\)，所以曲边梯形绕 \(x\) 轴旋转时，位置 \(x\) 处的截面是半径为 \(\sin x\) 的圆，面积为 \(\pi\sin^2x\). 由圆盘法，旋转体体积为
\[
V=\pi\int_0^\pi\sin^2x\,\mathrm{d}x
\]
利用半角恒等式 \(\sin^2x=\frac{1-\cos2x}{2}\)，有
\[
V=\pi\int_0^\pi\frac{1-\cos2x}{2}\,\mathrm{d}x
 =\pi\left[\frac{x}{2}-\frac{\sin2x}{4}\right]_{0}^{\pi}
 =\pi\left(\frac\pi2-0\right)=\frac{\pi^2}{2}
\]
所以旋转体的体积为 \(\frac{\pi^2}{2}\).
\end{solution}
